\section{Fourth physical variation and the exact high-order cancellation}
\label{sec:gr-fourth}
The third-variation result of Section~\ref{sec:el-localization} is retained.
It does not by itself bound a fourth derivative: differentiating the
stationary lift introduces its second derivative. The reflection
resolvent now controls that derivative without a small-incidence loss.
The following proof includes those terms rather than continuing the
third-order Schur cancellation beyond its valid scope.

Set
\begin{align}
 D_3&=R_++3R_+^2/\tau_-+3R_+^3/\tau_-^2,\\
 D_4&=R_++7R_+^2/\tau_-+18R_+^3/\tau_-^2+15R_+^4/\tau_-^3,\\
 \mathcal F_4&=2D_4C_1^4+6D_3C_2C_1^2,
 \qquad C_1=59/6,\quad C_2=2(a_*+a_*^2)C_1^2.
 \label{eq:gr-fourth-constant}
\end{align}
All constants concern the unchanged interval $R\in[9/20,47/100]$.

\begin{lemma}[Fourth tensor of one physical flight]\label{lem:gr-flight-four}
For $r=3,4$ and arbitrary collision-angle variations,
\[
 |D^r\tau_i[u^1,\ldots,u^r]|
       \le D_r\prod_{h=1}^r V_i(u^h),
 \qquad
 |D^rS[u^1,\ldots,u^r]|
       \le2D_r\prod_{h=1}^r\|u^h\|_1.
\]
\end{lemma}
\begin{proof}
For $N(X)=|X|$, put $v=X/|X|$ and $P=\mathrm{Id}-vv^{\mathsf T}$.
Then $DN=v^{\mathsf T}$, $D^2N=P/|X|$, and
\[
 D^3N[a,b,c]
 =-|X|^{-2}\{(v\cdot a)b^{\mathsf T}Pc
              +(v\cdot b)a^{\mathsf T}Pc
              +(v\cdot c)a^{\mathsf T}Pb\}.
\]
Thus $\|D^3N\|\le3|X|^{-2}$. Differentiating this display gives
$\|D^4N\|\le15|X|^{-3}$: differentiating the radial prefactor,
the three copies of $v$, and the three copies of $P$ contributes
at most $6$, $3$, and $6$, respectively, since
$Dv[d]=Pd/|X|$ and
$DP[d]=-\{(Pd)v^{\mathsf T}+v(Pd)^{\mathsf T}\}/|X|$.

Each derivative of a circular endpoint position has size at most
$R_+$ times the product of its angular variations. The fourth
composition rule has respectively $1,7,6,1$ set partitions into
$1,2,3,4$ blocks. Multiplying by the norm-derivative bounds
$1,|X|^{-1},3|X|^{-2},15|X|^{-3}$ proves the displayed $D_4$.
The third-order rule gives $D_3$, as in the retained proof.
Finally $\sum_i\prod_h V_i(u^h)\le2\prod_h\|u^h\|_1$,
by summing one factor and taking the suprema of the others.
This proves both full-action inequalities independently of $m$.
\end{proof}

\begin{theorem}[Collision-uniform fourth derivative of the reduced action]
\label{thm:gr-fourth}
On every actual regular endpoint branch and for all four endpoint
directions,
\begin{equation}\label{eq:gr-fourth-bound}
 |D_q^4F[\eta^1,\eta^2,\eta^3,\eta^4]|
       \le\mathcal F_4\prod_{h=1}^4\|\eta^h\|_\infty.
\end{equation}
There is no interior-incidence cutoff. In particular,
\[
 \|D_q^2G[\eta,\zeta]\|_{\rm op}
       \le2\mathcal F_4\|\eta\|_\infty\|\zeta\|_\infty.
\]
At fixed-itinerary multiple-grazing limits, these tensors have the
one-sided limits specified in Corollary~\ref{cor:gr-grazing-limit}.
\end{theorem}
\begin{proof}
Write $J_a=D_q\alpha[\eta^a]$ and
$K_{ab}=D_q^2\alpha[\eta^a,\eta^b]$. The latter has zero endpoint
coordinates. Differentiate the exact third-order identity
\eqref{eq:el-third-schur} in direction $\eta^4$. One obtains
\begin{equation}\label{eq:gr-fourth-identity}
 \begin{split}
 D_q^4F[\eta^1,\eta^2,\eta^3,\eta^4]
 ={}&D^4S[J_1,J_2,J_3,J_4]\\
 &+D^3S[K_{14},J_2,J_3]
  +D^3S[J_1,K_{24},J_3]
  +D^3S[J_1,J_2,K_{34}].
 \end{split}
\end{equation}
The three terms containing $K$ do not vanish. Their omission would
be an incorrect fourth-order use of stationary elimination.
Equivalently, let $g_{ab}=(D^3S[\mathbf e_i,J_a,J_b])_{i\in I}$.
Twice differentiated stationarity gives
$H_{II}(K_{ab})_I=-g_{ab}$. Hence the last three terms are
\[
 -g_{14}^{\mathsf T}H_{II}^{-1}g_{23}
 -g_{24}^{\mathsf T}H_{II}^{-1}g_{13}
 -g_{34}^{\mathsf T}H_{II}^{-1}g_{12}.
\]
This form verifies the symmetry but is not used to bound
$H_{II}^{-1}$ near grazing. The reflection equation gives instead
$\|J_a\|_1\le C_1\|\eta^a\|_\infty$ and
$\|K_{ab}\|_1\le C_2\|\eta^a\|_\infty\|\eta^b\|_\infty$.
Lemma~\ref{lem:gr-flight-four} applied to all four terms of
\eqref{eq:gr-fourth-identity} gives exactly
\eqref{eq:gr-fourth-constant}. For $m=1$, $K=0$ and the same bound
holds. Each entry of $D_q^2G[\eta,\zeta]$ is bounded by
$\mathcal F_4\|\eta\|_\infty\|\zeta\|_\infty$; a two-by-two
matrix has operator norm at most twice its largest entry.
The fixed-itinerary limit assertion follows from the smooth formal
reflection continuation already proved, restricted to the original
physical branch.
\end{proof}

\begin{corollary}[Uniform Taylor errors on the actual physical domain]
\label{cor:gr-taylor}
Suppose the entire straight segment $q+s h$, $0\le s\le1$, lies in
a single actual regular branch. Then
\begin{align}
 \left|F(q+h)-F(q)-DF(q)[h]-\tfrac12D^2F(q)[h,h]
                       -\tfrac16D^3F(q)[h,h,h]\right|
       &\le\tfrac1{24}\mathcal F_4\|h\|_\infty^4,\\
 \|G(q+h)-G(q)-DG(q)[h]\|_{\rm op}
       &\le\mathcal F_4\|h\|_\infty^2.
 \label{eq:gr-taylor}
\end{align}
The same bounds hold by a one-sided limit at a fixed-itinerary grazing
endpoint. They require neither a pressure estimate nor a prescribed
positive lower bound on intermediate incidences.
\end{corollary}
\begin{proof}
Apply the one-dimensional integral Taylor remainder to the stated
segment. The remainder integral for $F$ is
$\frac16\int_0^1(1-s)^3D^4F(q+sh)[h,h,h,h]\,ds$.
For $G$, integrate its second derivative with weight $1-s$.
The bounds follow from Theorem~\ref{thm:gr-fourth}; its one-sided
limit justifies the endpoint statement. The hypothesis that the
segment stays physical is not inferred from the derivative bound.
\end{proof}

\subsection{What this resolves in the original source problem}
The higher reflection jets and \eqref{eq:gr-fourth-identity} remove an
actual local analytic obstruction: repeated stationary differentiation
need not incur powers of the smallest visited incidence or an uncontrolled
factor equal to the collision count. They also justify the displayed
higher-order Taylor errors on every existing physical chart. These are
uniform statements about all actual regular itineraries, not assertions
from a finite sample or a replacement Markov model.

They do not count physical charts or continue through a newly hit disk.
The least unvisited clearance can still cut a chart, the original
section/return selector can still jump, and a positive one-sided paired
trace is not automatically a bulk-collision submeasure. In particular
no norm of the algebraic reflection resolvent is a pressure estimate
for the twisted transfer operator.

The remaining target is still \eqref{eq:v112-remaining}, with
$B\to\infty$ outermost and the full one-copy physical sum made before
the original radius, exact-label and roof suprema and collision limsup.
Neither \eqref{eq:gr-jets} nor \eqref{eq:gr-fourth-bound} supplies the
missing $m^{-2}$ common-physical-roof summation or nongrazing capacity.
The two averaged clauses of Input~\ref{input:v109-analytic}, the
stronger pressure/pairing interfaces and the original exact arithmetic
and conditioned-denominator hypotheses retain their distinct status.
No theorem or target is deleted to hide these remaining obligations.
